%% 第七章

\chapter{总结与展望}
\chaptermark{总结与展望}

\section{工作总结}

本文围绕低空无人机场景下多模态大模型的双认知时空推理能力评测与优化方法展开研究，整理并实现了 UAV-DualCog 双认知测试基准的本科论文初稿。与传统通用视觉问答或视频理解任务不同，无人机在开放三维空间中的观测具有具身性、动态性和多视角性。模型不仅需要理解外部目标、地标和环境变化，还需要理解无人机自身相对位置、未来视角变化和飞行行为。因此，本文将无人机具身推理拆分为自我状态认知和环境情况认知两条主线，并在图像和视频两种媒介下开展统一评测。

在测试基准设计方面，本文定义了六类任务。图像任务包括地标相对位置推理、未来观察预测、自我参照位置推理和地标驱动动作决策；视频任务包括飞行行为识别与时间定位、地标可见次数与区间推理。六类任务共同形成“能力维度、观测媒介、证据形式”的三维评测结构。本文还将目标框和时间区间作为正式评测对象，使实验能够区分语义正确性和证据可靠性。

在数据构建方面，本文实现了从仿真场景到测试样本的四阶段自动化工具链。该工具链包括场景扫描与语义点云融合、地标挖掘与人工复核、行为驱动视频任务构建、证据感知图像任务构建。当前测试基准覆盖 12 个测试场景、512 个有效地标、4096 个图像任务样本和 2048 个视频任务样本，并配套完成公开材料、网站展示、排行榜和分析页面整理。

在实验分析方面，本文对闭源模型、开源模型和空间推理微调模型进行了统一评测。结果表明，当前多模态大模型在 UAV-DualCog 上仍存在三类突出问题。第一，语义判断与显式证据之间存在脱节，模型可能答对选项却无法给出正确目标框或时间区间。第二，自我状态认知明显弱于环境情况认知，说明当前模型更容易理解外部目标，却更难建模无人机自身状态。第三，图像和视频之间存在真实但不充分的迁移，模型还没有形成稳定统一的跨媒介双认知能力。

在能力优化方法方面，本文根据评测发现完成了 UAV-DualCog-Train 训练集构建方案和两阶段模型优化路线的实现。训练集使用与测试集场景隔离的剩余场景构建；模型优化以 Qwen 3.5 4B 为基座，先通过监督微调学习“答案+证据”的结构化输出，再通过轻量 GRPO 引入格式、语义和空间证据奖励。与此同时，本文已经完成训练数据导出、完整性核查、训练脚本实现以及 smoke 和 tiny 规模验证。该实现把第五章中的能力短板转化为第六章中可执行、可验证的训练目标，使论文形成从基准构建、能力评估到能力优化验证的完整研究线路。

\section{主要结论}

本文可以得到以下结论。

第一，双认知能力是无人机具身推理的必要组成部分。对于无人机智能体而言，仅理解外部环境并不足够，模型还必须理解自身状态以及自身运动如何改变观测内容。UAV-DualCog 通过自我状态认知和环境情况认知两条主线，将这一能力要求显式纳入评测。

第二，语义正确性不能替代证据可靠性。实验中大量模型能够给出正确选项，却不能提供准确目标框或时间区间。这说明多模态大模型在无人机场景中的失败方式并非简单“看不懂”，而是常表现为语义判断和空间、时间证据之间的不一致。

第三，当前模型在双认知能力上存在显著失衡。所有配对模型的环境情况认知总体准确率均高于自我状态认知，总体均值差达到 16.97 个百分点。这表明自我状态建模仍是当前无人机多模态推理中的核心短板。

第四，模型规模和已有空间推理微调不能自动解决无人机场景问题。实验结果显示，更大规模模型、MoE 模型或一般空间推理微调模型并不必然在 UAV-DualCog 中取得更好结果。面向无人机双认知和结构化证据输出的任务内优化仍具有必要性。

第五，能力优化应围绕双认知协同和证据一致性展开。评测结果已经说明，单纯提高答案准确率不足以支撑可靠无人机智能体。当前已经实现的训练路线，正是同时围绕自我状态认知补强、环境情况认知保持、空间证据质量提升和跨媒介迁移验证展开，使模型逐步形成匹配、协同发展的双认知能力。

\section{后续展望}

后续工作将围绕第六章的正式能力优化实验继续推进。首先，需要等待当前正在进行的全量训练完成，并整理 SFT、GRPO Phase A 与 GRPO Phase B 的正式检查点和训练日志。其次，需要补充优化前后的定量结果表、双认知分支对比、视频迁移结果、消融实验和案例分析，重点检验格式奖励、语义奖励和空间证据奖励是否能够缩小语义与证据之间的差距。最后，需要结合正式结果修订第六章与本章总结，使毕业论文最终形成从测试基准构建、能力评估到能力优化验证的完整研究闭环。
